\section{Termín 3. února 2025}

\subsection{Tečná rovina}
Mějme rovinu $x^2yz + 3y^2 - 2xz^2 + 8z = 9$.
Nalezněte přímku kolmou k zadané rovině v bodě (1,-1,2) a její bod průniku s rovinou, která obsahuje souřadnicové osy 
$x$ a $z$.

\subsection{Použití diferenciálního počtu}
Určete hodnotu $3^{x+y+z}$ za předpokladu, že $x,y,z$ jsou kladná čísla a $\ln(xy^2z^3)$ má největší možnou hodnotu.

\subsection{Dvojný integrál}
Vypočítejte
\[
    \int_{0}^{4} \int_{\sqrt{x}}^{2} \frac{1}{y^3 + 1} \dif x \dif y
\]

\subsection{Potenciál vektorového pole}
Ověřte, že vektorové pole 
$\vec{F}(x, y, z) =
\begin{pmatrix}
    e^z + z^2 y, & \cos(y) + z^2 x, & xe^z + 2xyz
\end{pmatrix}$ je konzervativní a nalezněte jeho potenciál.

\subsection{Plošný integrál}
Spočtěte plošný integrál $\iint\limits_{M} \vec{F} \dif \vec{s}$, kde 
\[\vec{F}(x,y,z) = 
    \begin{pmatrix}
        -xz, -yz, z^2 + z
    \end{pmatrix}
\] 
a $M$ je povrch elipsoidu 
\[
    \frac{x^2}{a^2} + \frac{y^2}{b^2} + \frac{z^2}{c^2} =1, a,b,c > 0 \text{ s vnější orientací.}
\]